%%%%%%%%%%%%Outline%%%%%%%%%%%%%%%%%%

%%%
% message<the eBPF machine-code gap from a small, verifier-checkable ISA>
% eBPF widely deployed for networking, observability, security on hot paths
% yet, for one source, eBPF's machine code slower than direct native compilation
% the gap comes from an instruction set omitting a native compiler's hardware instructions
% the instruction set stays small so the verifier can check every program

%%%
% message<\tool, a verified extension mechanism for the eBPF compilation pipeline>
% \tool extends the eBPF compilation pipeline with verified extensions
% each extension has a verifier-checked bytecode form and a JIT-emitted native form
% a proof equates the two forms, so the verifier alone secures safety
% \lang instantiates \tool as extensions for hardware instruction idioms
% eval: XX faster than the existing eBPF JIT on x86-64 and ARM64, small kernel change

%%%%%%%%%%%%Outline%%%%%%%%%%%%%%%%%%

% [to-do] zhengjie: switch template
% [to-do] zhengjie: check terminalogy
% - kernel JIT -> eBPF JIT compiler
% - BPF - eBPF
% - eBPF compilation pipeline with \tool

\begin{abstract}
eBPF is widely deployed for networking, observability, and security, attaching to hot kernel paths. However, for the same source program, the eBPF pipeline produces machine code that runs slower than a native compiler's machine code. The gap comes from the eBPF instruction set, which omits the hardware instructions a native compiler selects for speed. The eBPF instruction set stays small so that the verifier can check every program before execution.
%
We present \tool, which extends the eBPF compilation pipeline with verified extensions. Each \tool extension provides two forms of the same operation: a bytecode form that the existing verifier checks, and a native form that the JIT compiler emits. Because a proof guarantees that the native form behaves like the bytecode form, the existing verifier alone secures safety. We instantiate \tool as \lang, a set of extensions for hardware instruction idioms such as rotation, conditional select, and bitfield extraction. On x86-64 and ARM64, \lang makes eBPF programs run XX\% faster than the existing eBPF JIT compiler with a small kernel change.
\end{abstract}
